\section{Introduction}
Frank--Wolfe (FW) methods solve linear subproblems in place of
projections. Backtracking makes their quadratic step models usable
without a prescribed smoothness constant, and away steps improve
progress on polytopes~\cite{braun,pedregosa,lacostejaggi}.
A global curvature bound can nevertheless discard substantial
information when curvature decreases with the objective value.
We consider the affine loss-suboptimality Hessian condition studied in
\cite{warmup}, specialized to a constrained convex problem. It is also
the exponent-one case of the loss-dependent smoothness family in
\cite{liuwarmup}. More general
function-value bounds on Hessians already appear in self-bounding
regularity frameworks~\cite{kale,caoselfbound}; the condition itself is
not a contribution of this paper.

Our question is quantitative: what does ordinary backtracking guarantee
for an observable stopping rule under this condition, including its
initial estimate and all trial evaluations? Theorem~\ref{thm:fw}
answers it with separate $H_0/\varepsilon$ and $H_1\log(1/\varepsilon)$
contributions. The main device is a potential for estimates that remain
too large as the certificate enters lower dyadic levels. It does not
change the algorithm or restrict the LMO to a sublevel set.
Theorem~\ref{thm:afw} extends the result by standard atom-removal
counting. Under a geometric gap inequality,
Theorem~\ref{thm:geometric} gives logarithmic accuracy dependence with
an explicit initial $H_1$ cost. Theorem~\ref{lem:bridge} gives the
optimal pointwise FW-gap envelope for this class, with attaining
one-dimensional objectives. Proposition~\ref{prop:zeroface} characterizes
the zero-base case $H_0=0$: its entire solution set is an exposed face
with a common gradient. Theorem~\ref{thm:zerofinite} then upgrades
approximate stopping to finite exact termination of backtracking AFW,
with an explicit bound independent of requested accuracy. This uses
known identification mechanisms~\cite{bomze,garber}; neither a new AFW
algorithm nor general identification theory is asserted.
Complete proofs, boundary examples and counterexamples are provided.

\paragraph{Related analyses.}
Function-dependent smoothness and two-regime gradient-method bounds
precede this work~\cite{warmup,lobanovkoloskova,lobanov}.
Vyguzov and Stonyakin~\cite{vyguzov} analyze FW under gradient-norm
dependent smoothness. On a compact convex domain,
$f-f_Q^*\le D\norm{\nabla f}$ converts our condition to one such Hessian
bound, but a global bound on the largest gradient does not retain the
separated dependence studied here.
Pedregosa et al.~\cite[Theorems 1, 3, and 4]{pedregosa} already give
backtracking trial counts, an averaged-gradient Fenchel dual guarantee,
and linear AFW convergence with accepted local curvature estimates.
Our scalar lower-bound maximum differs from that Fenchel dual; their
values cannot simply be identified. The result here specializes local
curvature to a certificate threshold and explicitly accounts for lag.
It does not claim priority for backtracking, dual certificates, linear
AFW convergence, or the telescoping identity for rejected trials.
Best-so-far tangent bounds and their computable stopping gaps are
also classical~\cite[Method 1]{freundgrigas}. Optimistic FW and
primal--dual short steps explicitly optimize computable gaps under
global smoothness~\cite[Proposition 4.1]{beyondshort}; they do not
evaluate the inherited halving/doubling rule analyzed here.

Other related frameworks include generalized self-concordant FW,
adaptive Bregman steps, and adaptive conditional gradient
sliding~\cite{selfconcordant,bregmanfw,adcgs}.
Extended-curvature bounds and restarted error-bound schemes already
provide two-regime and geometric rates~\cite{wirth,restartfw};
linear objective- and FW-gap convergence beyond global gradient
smoothness is also known for log-homogeneous barriers~\cite{zhao}.
Auto-conditioned FW supplies a practical comparator with different
preprocessing and oracle accounting~\cite{autoconditioned}.
Gradient-difference initialization with backtracking~\cite{khademi}
and line-search-free conditional gradients~\cite{yagishita}
provide further parameter-adaptive alternatives.
ALFCG estimates smoothness without line search but retains global
component or average smoothness assumptions and a stochastic
stationarity criterion~\cite[Section 3]{alfcg}.
Recent damped Newton--FW uses Hessian information and assumes positive
definite Hessians~\cite[Section 1]{newtonfw}; this does not cover our
zero-base optimum, where the Hessian vanishes.
Appendix~\ref{app:comparison} specifies the differences used in our
comparison. A direct global-smoothness substitution is generally weaker,
but that observation is not a claim that no equivalent consequence
can be obtained from another adaptive analysis.

The empirical study addresses three issues relevant to interpreting
these bounds: whether certificate initialization helps independently of
a starting-point change; whether away steps resolve active-face
obstructions; and whether failure to stop reflects remaining primal
error or a conservative certificate. Oracle budgets are kept separate,
and failed runs are retained.
